\begin{enumerate}
    \item Nuskaitomas duomenų failas eilutė po eilutės. Eilutės, kuriose yra trūkstamų reikšmių (pažymėtų „?“), yra praleidžiamos.
    \item Pašalinamas pirmasis stulpelis (paciento ID), paliekant tik $9$ požymių stulpelius ir klasės stulpelį.
    \item Klasės reikšmės konvertuojamos:
        \begin{itemize}
            \item $2 \rightarrow 0$ (nepiktybinis).
            \item $4 \rightarrow 1$ (piktybinis).
        \end{itemize}
    \item Duomenų eilutės permaišomos atsitiktine tvarka naudojant fiksuotą atsitiktinių skaičių generatoriaus pradinę reikšmę ($seed=42$), kad rezultatai būtų atkartojami.
    \item Požymių reikšmės normalizuojamos į intervalą $[0; 1]$ naudojant min-max~\eqref{eq:min-max} metodą pagal formulę:
    $$x_{norm} = \frac{x - x_{min}}{x_{max} - x_{min}}$$
    \item Metodas \texttt{save()} sujungia normalizuotus požymius su klasės stulpeliu ir išsaugo rezultatą į naują CSV failą.
\end{enumerate}